\chapter{Especificación técnica}

\section{Problema seleccionado}

Asignación de salones y franjas horarias a bloques académicos. Cada bloque necesita un salón compatible con su tipo y no puede coincidir con otros bloques que compartan estudiantes, profesor o franja horaria.

\section{Cómo se modela en el código}

El proyecto representa cada bloque como un vértice y cada conflicto como una arista. Los salones disponibles se usan como identificadores de asignación. El constructor del grafo recibe materias, grupos y salones, y devuelve la estructura que después consume el solucionador de horarios.

La implementación que genera el horario global se apoya en tres piezas:

\begin{itemize}[leftmargin=1.5cm]
  \item el cargador de datos, que normaliza el JSON de entrada;
  \item el constructor del grafo de conflictos, que detecta incompatibilidades;
  \item el solucionador global, que asigna salones y produce el horario semanal.
\end{itemize}

\section{Algoritmos seleccionados}

El proyecto usa dos familias de algoritmos en la ruta principal:

\begin{itemize}[leftmargin=1.5cm]
  \item coloreo de grafos para la asignación global de salones;
  \item Dijkstra para construir horarios personalizados de estudiantes cuando se necesitan rutas de menor costo entre clases.
\end{itemize}

\section{Entrada y salida}

Entrada: archivos JSON con materias, salones y franjas. Salida: horario semanal global y horarios individuales para estudiantes, ambos en estructuras JSON/TypeScript que pueden consumirse desde la API.

\section{Validación básica}

Con los datos incluidos en el repositorio se puede generar un horario global y recuperar la información del grafo. Para horarios estudiantiles, la validación consiste en revisar que el conjunto de materias solicitado tenga una asignación compatible y que no exista solapamiento entre bloques.
